\documentclass{article}

\usepackage[utf8]{inputenc}  
\usepackage{graphicx} 
\usepackage{amsmath} 
\usepackage[version=4]{mhchem} 
\usepackage{siunitx} 
\usepackage{longtable,tabularx}   
\usepackage{amssymb, bm} % Removed duplicate amsmath load   
\usepackage{verbatim}    
\usepackage{listings}  
\usepackage{xcolor} % Added missing package for \definecolor

\definecolor{codegray}{rgb}{0.5,0.5,0.5} 
\definecolor{codepurple}{rgb}{0.58,0,0.82} 
\definecolor{backcolour}{rgb}{0.95,0.95,0.92}

\lstdefinestyle{pythonstyle}{
    backgroundcolor=\color{backcolour},
    commentstyle=\color{codegray},
    keywordstyle=\color{blue},
    numberstyle=\tiny\color{codegray},
    stringstyle=\color{codepurple},
    basicstyle=\ttfamily\footnotesize,
    breakatwhitespace=false,         
    breaklines=true,                 
    captionpos=b,                    
    keepspaces=true,                 
    numbers=left,                    
    numbersep=5pt,                  
    showspaces=false,                
    showstringspaces=false,
    showtabs=false,                  
    tabsize=4,
    language=Python
}

\lstset{style=pythonstyle}


\setlength\LTleft{0pt} 



%%%%%%%%%%%%%%%%%%%%%%%%%% HELPFUL MATH DEFINITIONS %%%%%%%%%%%%%%%%%%%%%%%%%%%%%%

% vectors / matrices
\newcommand{\vect}[1]{\mathbf{#1}}
\newcommand{\mat}[1]{\mathbf{#1}}

% position / velocity / quaternion shorthands
\newcommand{\rBM}{\vect{r}^M_{B/M}}      % position of B wrt M expressed in M
\newcommand{\vBM}{\vect{v}^M_{B/M}}      % velocity of B wrt M expressed in M
\newcommand{\qBM}{\hat{\Bar{\mathbf{q}}}^B_{M}}    % quaternion representing B wrt M (unit)

\newcommand{\mrBM}{\vect{m}_{\rBM}^{-}}      % position of B wrt M expressed in M

\newcommand{\mrBMpkm}{\vect{m}_{r,k-1}^+}

% MEKF defs
\newcommand{\qUnitGeneric}{\hat{\Bar{\mathbf{q}}}} 
\newcommand{\qRef}{^*\hat{\Bar{\mathbf{q}}}^B_{M}}
\newcommand{\wBMRef}{^*\vect{\omega}^B_{B/M}}
\newcommand{\errQuat}{\delta\qUnitGeneric^{B}_{B'}}
\newcommand{\errQuatInv}{{\delta\qUnitGeneric^{B}_{B'}}^{-1}}
\newcommand{\errwBM}{\delta\vect{\omega}^B_{B/M}}

% MEKF gryo
\newcommand{\wBMMeas}{\Tilde{\vect{\omega}}^B_{B/M}}
\newcommand{\wBias}{\vect{\beta}}
\newcommand{\wNoise}{\vect{w}_\omega}
\newcommand{\wBiasNoise}{\vect{w}_\beta}
\newcommand{\PSDww}{\mat{Q}_\omega\omega}
\newcommand{\PSDbb}{\mat{Q}_\omega\omega}

\newcommand{\mGyro}{\vect{m}_\omega}
\newcommand{\mBias}{\vect{m}_\beta}


\newcommand{\TBMprime}{\mat{T}_{M_k}^{B'}}




% dynamics variables
\newcommand{\wMN}{\vect{\omega}^M_{M/N}}   
\newcommand{\muMoon}{\mu_{Moon}}   
\newcommand{\rBMnorm}{r^M_{B/M}}  
\newcommand{\wBM}{\vect{\omega}^B_{B/M}}
\newcommand{\wGryoOut}{\vect{\omega}^G_{B/N}}

% measurement variables
\newcommand{\rlO}{\vect{r}^M_{l/M}}
\newcommand{\rBO}{\vect{r}^M_{B/M_,k}}
\newcommand{\rlB}{\vect{r}^M_{l/B_,k}}
\newcommand{\rlBmeas}{\vect{r}^C_{l/C_,k}}
\newcommand{\TBM}{\mat{T}^B_{M_,k}} 
\newcommand{\TCB}{\mat{T}^C_{B}} 

%% ^^ above is old code, but has the notation that I use

\newcommand{\pBiB}{\vect{p}^{\mathit{B}}_{q_{i}/\mathbf{B}}} 
\newcommand{\pBjB}{\vect{p}^{\mathit{B}}_{q_{j}/\mathbf{B}}} 
\newcommand{\pBi}{\vect{p}_{q_{i}/\mathbf{B}}}
\newcommand{\pBip}{\vect{p}_{p_{i}/\mathbf{B}}} 
\newcommand{\pCiC}{\vect{p}^{\mathit{C}}_{p_{i}/\mathbf{C}}} 
\newcommand{\pCi}{\vect{p}_{p_{i}/\mathbf{C}}} 
\newcommand{\pBCC}{\vect{p}^{\mathit{C}}_{\mathbf{B}/\mathbf{C}}} 
\newcommand{\pBC}{\vect{p}_{\mathbf{B}/\mathbf{C}}} 
\newcommand{\TBCtoC}{\mat{T}_{\mathit{B} \rightarrow \mathit{C}}} % 4x4 Homogeneous Transformation 
\newcommand{\nCiC}{\vect{n}^{\mathit{C}}_{p_{i}/\mathbf{C}}} % Surface normal vector at target point p_i in C



%% RSO is point in 3d space
%% measurement model that 

\newcommand{\TBC}{\mat{T}^B_{C}} 
\newcommand{\TIB}{\mat{T}^I_{B}} 
\newcommand{\uRSOcC}{\vect{u}^C_{RSO/C}}
\newcommand{\uRSOcI}{\vect{u}^I_{RSO/C}}

\begin{document}

\section{3D point cloud matching geometry:}

To transform points on a rigid body frame $B$ to points in a camera frame $C$, we can express the operation using both a 3D affine vector notation and a $4 \times 4$ homogeneous transformation matrix.

\subsection{Frame Definitions}

\begin{equation}
   \mathit{C} = \{ \mathbf{C}, \hat{\vect{c}}_1, \hat{\vect{c}}_2, \hat{\vect{c}}_3 \}
\end{equation}
    
\begin{equation}
   \mathit{B} = \{ \mathbf{B}, \hat{\vect{b}}_1, \hat{\vect{b}}_2, \hat{\vect{b}}_3 \}
\end{equation}

\subsection{Measurement Point Cloud Point Defintion}
The point cloud derived from measurement source has $P$ number of points, which mathematically are vectors in $\mathbb{R}^3$ of each point wrt the camera origin ,$\mathbf{C}$ , expressed in the $\mathit{C}$ frame.

\begin{equation}
\pCiC , i = 1, 2, 3,... P
\end{equation}

\subsection{CAD Model Point Cloud Points Defintion}
The point cloud derived from the CAD model may not have the same number of points as the measurement itself. The CAD model point cloud has $Q$ number of points, which mathematically are vectors in $\mathbb{R}^3$ of each point wrt the CAD PCD origin ,$\mathbf{B}$ , expressed in the $\mathit{B}$ frame.

\begin{equation}
\pBjB , j = 1, 2, 3,... Q
\end{equation}

\subsection{Measurement Model As The Transformation in $\mathbb{R}^3$}
The standard 3D kinematic translation and rotation mapping is given by:
\begin{equation}
    \pCiC = \TCB \pBjB + \pBCC
\end{equation}

\subsection{The Coorespondence Problem}
Each measurement $\pCiC$ has an assocated $\pBjB$ such that in a perfect noiseless environment, $\TCB$ and $\pBCC$ would map each $\pBjB$ to $\pCiC$ perfectly.

Therefore we define the coorespondence set 
\begin{equation}
K = \{ (p_{i},q_{j}) \}
\end{equation}


Note that in the above equation point, even in a noiseless environment $p_{i}$ and $q_{j}$ are not gaurenteed to be the same point due to the discretizate sampling of the CAD model itself.


\subsection{Generic Homogeneous Transformation in $\mathbb{R}^4$}
By appending a scaling factor of 1 to our coordinate vectors, the rigid body mapping can be cleanly expressed as a single $4 \times 4$ matrix multiplication:
\begin{equation}
    \begin{bmatrix}
        \pCiC \\
        1
    \end{bmatrix}
    = \TBCtoC
    \begin{bmatrix}
        \pBiB \\
        1
    \end{bmatrix}
\end{equation}

Where the homogeneous transformation matrix $\TBCtoC$ is explicitly partitioned as:
\begin{equation}
    \TBCtoC = 
    \begin{bmatrix}
        \TCB & \pBCC \\
        \vect{0}_{1\times3} & 1
    \end{bmatrix}
\end{equation}

Expanding this reveals the complete linear algebraic formulation:
\begin{equation}
    \begin{bmatrix}
        \pCiC \\
        1
    \end{bmatrix}
    =
    \begin{bmatrix}
        & & &  \\
        & \TCB & & \pBCC \\
        & & &  \\
        &\vect{0}^T & & 1
    \end{bmatrix}
    \begin{bmatrix}
        \pBiB \\
        1
    \end{bmatrix}
\end{equation}


\subsection{Point-to-Plane Cost Function}
The point-to-plane error metric minimizes the squared distance along the surface normal of the target point in the camera frame $C$. 

Let $\pBjB$ represent a source point in the rigid body frame $B$, $\pCiC$ represent its corresponding target point on the surface in the camera frame $C$, and $\nCiC$ be the unit surface normal vector at $\pCiC$.

The objective function seeking the optimal rigid transformation matrix $\TBCtoC$ is formulated as:

\begin{equation}
    J = \sum_{K = \{ (p_{i},q_{j}) \}} \left[ \left( \pCiC - (\TCB \pBjB + \pBCC) \right) \cdot \nCiC \right]^2
\end{equation}


\begin{equation}
    J = \sum_{K = \{ (p_{i},q_{j}) \}} \left[ \left( \pCiC - (\TCB \pBjB + \pBCC) \right) \right]^2
\end{equation}

\begin{equation}
    \min_{\TCB, \pBCC} (J)
\end{equation}

Expressed compactly in homogeneous coordinates ($\mathbb{R}^4$), let $\tilde{\vect{p}}$ denote the homogeneous vector $[ \vect{p}^T, 1 ]^T$. The cost function is written as:

\begin{equation}
    E(\TBCtoC) = \sum_{i} \left[ \left( \mathbf{\Pi} \TBCtoC \tilde{\vect{p}}^B_{q_{i}/B} - \pCiC \right)^T \nCiC \right]^2
\end{equation}

Where $\mathbf{\Pi} = \begin{bmatrix} \mat{I}_{3\times3} & \vect{0}_{3\times1} \end{bmatrix}$ is a projection matrix used to drop the homogeneous scaling factor back down to $\mathbb{R}^3$ for the dot product calculation.

\subsection{Laying out some uncoordinatized vectors}

\begin{equation}
    \pBi
\end{equation}
\begin{equation}
    \pCi
\end{equation}
\begin{equation}
    \pBC
\end{equation}
\begin{equation}
    \pBip
\end{equation}


\newpage
\subsection{RSO Direction Transformation Measurement Model}
The unit line-of-sight vector to the Resident Space Object (RSO) observed in the camera frame $C$, denoted as $\uRSOcC$, can be mapped to the inertial frame $I$ by chaining the rigid body rotation transformations. 

The complete geometric measurement model is written as:
\begin{equation}
    \uRSOcI = \TIB \TBC \uRSOcC
\end{equation}

\begin{equation}
   \mathit{N} = \{ \mathbf{O}, \hat{\vect{n}}_1, \hat{\vect{n}}_2, \hat{\vect{n}}_3 \}
\end{equation}
\begin{equation}
    \mathit{I} = \{ \mathbf{O}, \hat{\vect{i}}_1, \hat{\vect{i}}_2, \hat{\vect{i}}_3 \}
\end{equation}
\begin{equation}
   \mathit{M} = \{ \mathbf{M}, \hat{\vect{m}}_1, \hat{\vect{m}}_2, \hat{\vect{m}}_3 \}
\end{equation}

\end{document}